% !TEX root = ../main.tex

\chapter{绪论}

\section{研究背景与意义}

鞋履设计是兼具功能性、审美性和制造约束的产品设计任务。设计师在早期方案阶段通常会使用侧视图、俯视图、底视图等二维草图表达鞋面轮廓、鞋底结构、鞋带区域、后跟支撑、装饰线条和比例关系。二维草图具有绘制成本低、修改灵活和表达直接的优点，但它不能充分呈现鞋履在空间中的体积关系，也难以直观检查鞋楦包覆、鞋底厚度、侧墙转折和内部空腔等三维结构。若要从草图进入可旋转预览或三维资产制作阶段，传统流程往往需要建模人员在三维软件中手工建模、修面和调整材质，成本较高，迭代周期较长。

近年来，深度学习方法推动了图像到三维模型生成的发展。早期工作尝试从单张图像直接预测三维网格结构，例如Pixel2Mesh和Mesh R-CNN等方法将二维视觉特征与三维形状预测相结合\citep{wang2018pixel2mesh,gkioxari2019meshrcnn}。随后，NeRF等神经场表示方法提升了新视角合成和隐式场景表达能力\citep{mildenhall2020nerf}。扩散模型进一步带来了利用二维生成先验进行三维生成的路线，例如DreamFusion利用文本到图像扩散先验优化三维表示\citep{poole2022dreamfusion}，Zero-1-to-3、One-2-3-45、Wonder3D等方法则围绕单图像新视角生成和三维重建展开\citep{liu2023zero123,liu2023one2345,long2023wonder3d}。这些工作表明，从二维输入快速生成三维资产已经具备较强的技术可行性。

然而，鞋履设计草图到三维模型的重建仍然具有明显挑战。首先，设计草图与自然图像或产品渲染图存在显著域差异。草图通常由线条、少量明暗和设计标注组成，缺少纹理、材质和真实光照信息，而多数图像到三维模型更适合处理自然图像或渲染图。其次，鞋履本身具有薄壁、空腔、鞋口、鞋底层次和局部细节等复杂结构，仅凭单视图很难推断不可见的内部几何。再次，毕业设计演示需要在本地环境中形成可复现流程，这对模型选择、推理参数、文件导出和交互界面都提出了工程约束。因此，如何在可控的本地软硬件条件下，构建一个面向鞋履草图输入的可运行三维重建原型系统，具有实际的研究和工程意义。

本文围绕“融合三视图特征的鞋履3D模型重建研究”这一题目，关注从单视图或三视图鞋履设计草图生成可预览三维鞋履模型的流程。本文不将目标设定为直接生成可制造级CAD/NURBS模型，而是以可导出的GLB网格模型、可视化演示和可复现实验分析作为阶段性目标。这样的定位更符合当前开源图像到三维模型的能力边界，也便于在论文中清晰区分可实现的外观重建能力和仍未解决的内部结构推断问题。

\section{问题定义}

本文研究的问题可以表述为：给定一张单视图鞋履设计草图，或由侧视图、俯视图、底视图等组成的三视图草图输入，系统需要经过图像预处理、草图域适配、三维形状生成和网格质量分析等步骤，输出一个可在浏览器或三维查看器中预览的鞋履GLB模型，并给出与论文实验分析相关的网格指标和风险提示。

从输入角度看，本文关注的正式输入是鞋履设计草图，而不是普通商品照片。商品照片可以作为模型可用性的早期验证样例，但它并不能代表最终任务。草图输入通常包含较强的结构意图，例如鞋底轮廓、鞋面分割线、鞋带区域和后跟形状；同时也缺少真实图像中的连续色彩、材质和阴影。为了降低这种域差异，本文将ControlNet作为草图到模型友好图像之间的适配模块。ControlNet能够在扩散模型生成过程中引入边缘、草图或其他条件控制信号\citep{zhang2023controlnet}，因此适合用于将鞋履线稿转换为更接近产品渲染图的中间输入。

从输出角度看，本文输出的是网格资产而不是参数化CAD模型。输出模型需要满足三个基本要求：第一，文件能够成功导出并被常见三维查看器或Gradio界面加载；第二，外部轮廓应尽可能保留鞋履输入中的鞋头、鞋底、鞋面、后跟等主要设计信息；第三，系统需要对网格的面数、包围盒、厚度代理指标、连通组件和内部结构不确定性进行记录。特别需要指出的是，鞋履内部结构属于单视图或有限视图条件下的隐藏几何推断问题。当前系统可以生成外观上较为合理的鞋履模型，但不能保证内部空腔、鞋口厚度和真实穿着空间完全正确，因此本文将其作为实验分析和未来工作的重点讨论对象。

从工程约束角度看，本文系统需要在本地WSL环境下运行，并兼顾模型可部署性、GLB导出稳定性、交互演示需求和后续扩展空间。本文没有把模型路线选择写成单一模型的离线比较，而是围绕“能否服务鞋履草图到三维预览”这一目标进行工程取舍。SF3D能够在本地生成可读取的GLB模型，可作为稳定基线\citep{boss2024sf3d}；Hunyuan3D-2mini在本项目鞋履样例上表现出更好的外观质量和较好的本地可行性，因此被选为当前主要形状生成后端\citep{hunyuan3d2025hunyuan3d2}。对于三视图输入，本文在数据组织和方法设计上保留侧视图、俯视图和底视图的接口，并将其作为后续优化主路线的重要方向。

\section{本文主要工作}

围绕上述问题，本文完成的主要工作如下。

第一，本文调研并筛选了适合鞋履草图任务的开源图像到三维模型路线。模型选择过程综合考虑GLB导出能力、推理稳定性、部署复杂度、鞋履样例视觉效果和后续三视图扩展潜力，对系统后端进行了重新定位。最终系统以Hunyuan3D-2mini作为当前默认形状生成后端，以SF3D作为稳定基线和对照方案。

第二，本文构建了面向鞋履设计草图的预处理和草图域适配流程。系统能够对输入草图进行尺寸归一化、背景整理和边缘控制图生成，并通过ControlNet将线稿式输入转换为更适合图像到三维模型处理的中间渲染图。该流程保留了原始草图、控制图和生成图，便于后续论文截图、质量对比和局限性分析。

第三，本文实现了端到端的鞋履三维重建原型系统。系统包含命令行流水线和Gradio演示界面，支持选择直接输入、ControlNet渲染输入或二者对比，支持Hunyuan3D和SF3D后端，并能够输出GLB模型、关键中间图像、运行摘要和网格报告。该系统使从草图输入到三维模型预览的流程具备可操作性和可展示性。

第四，本文设计了非破坏式网格质量报告模块。该模块不直接修改生成网格，而是记录GLB加载状态、文件大小、几何数量、顶点数、面数、包围盒、尺寸比例、连通组件和水密性等指标，并对薄片化、碎片化和内部结构不可验证等风险给出提示。该报告既服务于演示界面，也为论文实验章节提供可复现的定量依据。

第五，本文对系统局限进行了明确界定。当前系统能够在本地硬件上完成鞋履外观级三维模型生成，但对真实内部结构、制造级曲面质量和严格三视图一致性仍不能保证。本文将这些问题作为后续微调、鞋履草图数据集建设、多视图约束和结构先验建模的研究方向，而不是在当前阶段过度宣称已经解决。

\clearpage

\section{论文组织结构}

本文共分为六章。

第一章为绪论，介绍鞋履草图到三维模型重建的研究背景、任务定义、主要工作和论文结构。第二章为相关工作，回顾单图像三维重建、扩散先验三维生成、多视图生成、前馈式三维重建以及草图条件控制等研究方向。第三章为方法设计，说明本文系统的总体流程、草图输入处理、ControlNet适配、SF3D基线、Hunyuan3D生成路线和网格质量报告设计。第四章为系统设计与实现，介绍项目环境、模块划分、命令行流水线、后端封装和Gradio演示界面。第五章为实验与分析，先展示SF3D基线结果，再分析Hunyuan3D主路线表现，并讨论网格指标、视觉效果、三视图融合和内部结构问题。第六章为总结与展望，总结本文完成的工作，分析当前系统局限，并提出未来在数据集、模型微调、多视图融合和结构约束方面的改进方向。
